% Style note (v2, 2026-05-07): rewritten in the EMNLP 2025 Outstanding-paper
% style — argument-first, no numeric headline figures, verbs of inquiry. The
% specific call counts and matrix sizes live in §1 and §3, not the abstract.

Setwise large language model (LLM) rerankers typically exploit a small candidate window: each LLM call sees only a few documents, and a sorting procedure over many such calls produces the final ranking. This locality was a context-length compromise. We ask whether it is still necessary, given that current open-weight long-context LLMs can fit an entire first-stage reranking pool into a single prompt. We introduce \textbf{MaxContext}, a zero-shot setwise reranker family in which every call sees the whole live pool. MaxContext spans three output policies: best-only selection (\textsc{TopDown}), worst-only selection (\textsc{BottomUp}), and joint best--worst selection (\textsc{DualEnd}). The three variants share an interface but differ deterministically in serial LLM calls per query: \textsc{DualEnd} removes two documents per call, while either single-extreme variant removes one. We argue that this design isolates two questions that are usually conflated --- whether whole-pool selection is competitive with windowed setwise sorting, and whether eliciting both ends of the ranking jointly preserves quality while approximately halving serial LLM calls --- and we test both across nine open-weight rerankers from the Qwen3.5, Llama-3.1, and Ministral-3 families on TREC DL19/20 and six BEIR domains. Our analyses also surface input-order effects under whole-pool prompting that small-window setwise reranking does not expose, suggesting that the long-context regime is not order-invariant.
